% chapters/ch36_appendix_c_toys.tex
% Chapter 36 (Part VI): Appendix C — Illustrative Toy Simulations (Non-normative)
% Source: NRMO_v5_updated.pdf, Ch36 Appendix C (lines 13000-13129 of v5_layout.txt)

\chapter{Appendix C --- Illustrative Toy Simulations
(Non-Normative)}
\label{ch:appendix-c-toys}

\nrmobox{Document role}{
\textbf{Status}: Non-normative illustrative implementation.\\[0.3em]
This appendix presents two toy simulations demonstrating NRMO
behaviour: a buffer-policy comparison (C.1) and a safety-certificate
construction with endogenous / exogenous hazard separation (C.2).
Code is illustrative; the canonical implementation is in Part~X.
}

% =====================================================================
\section{Appendix C.1 --- Buffer Policy Comparison}
\label{sec:appendix-c-1}

This first toy illustrates the comparison between unconstrained
greedy, buffer-modulated, and NRMO-core-with-venture policies. The
canonical implementation pattern:

\begin{lstlisting}[language=Python,basicstyle=\ttfamily\scriptsize,
                   frame=single]
def policy_buffer(env0, B_threshold=0.30):
    def pi(s, B):
        x, y = s
        if B < B_threshold:
            # Buffer low: defensive
            if y < 4:
                return 2  # move down (away from cliff)
            return 0
        # Buffer adequate: progress toward goal
        if y < 4:
            return 2
        if x < 4:
            return 1
        return 0
    return pi

def policy_nrmo_venture(env0, venture_p=0.05):
    def pi(s, B):
        x, y = s
        near_cliff = any(
            (x + dx, y + dy) in env0.cliff
            for dx, dy in [(0, 1), (1, 0), (-1, 0), (0, -1)]
        )
        if (not near_cliff) and random.random() < venture_p:
            return 1
        if y < 4:
            return 2
        if x < 4:
            return 1
        return 0
    return pi

# Results (10,000 episodes, slip=0.10):
#   Unconstrained greedy:    ruin_rate = 0.0842
#   Buffer policy:           ruin_rate = 0.0021
#   NRMO core + venture 5%:  ruin_rate = 0.0031
\end{lstlisting}

\paragraph{Interpretation.}
\begin{itemize}
\item Unconstrained greedy: ruin rate 8.42\%.
\item Buffer policy: ruin rate 0.21\% --- $\sim 40\times$ reduction.
\item NRMO core + 5\% venture: ruin rate 0.31\% --- maintains low
ruin while permitting bounded exploration.
\end{itemize}

% =====================================================================
\section{Appendix C.2 --- Safety Certificate
(Endogenous vs Exogenous Hazards)}
\label{sec:appendix-c-2}

This second toy illustrates:

\begin{enumerate}
\item separating endogenous failure from exogenous background
hazard;
\item irreversibility proxy $I(s, a)$;
\item bounded-horizon certificate $\Pr(\tau_F \le T) \le \varepsilon$
estimated by Monte Carlo.
\end{enumerate}

\begin{lstlisting}[language=Python,basicstyle=\ttfamily\scriptsize,
                   frame=single]
import random, math
from dataclasses import dataclass

@dataclass
class ToyWorld:
    w: int = 6
    h: int = 4
    slip:  float = 0.10
    goal       = (5, 0)
    start      = (0, 3)
    p_exogenous: float = 1e-4
    F_endogenous = frozenset({(2, 1), (3, 1), (4, 1)})

def irreversibility_I(world, s, a, kappa=1.0):
    def dist_F(xy):
        return min(abs(xy[0] - fx) + abs(xy[1] - fy)
                   for (fx, fy) in world.F_endogenous)
    dx, dy = [(0, -1), (1, 0), (0, 1), (-1, 0)][a]
    ns = (max(0, min(world.w - 1, s[0] + dx)),
          max(0, min(world.h - 1, s[1] + dy)))
    d0, d1 = dist_F(s), dist_F(ns)
    penalty = 2.0 if d1 == 0 else (1.0 if d1 == 1 else 0.0)
    J = max(0.0, float(d0 - d1)) + penalty
    return 1.0 - math.exp(-kappa * J)

def policy_nrmo_like(world, s, I_threshold=0.55):
    safe = [a for a in range(4)
            if irreversibility_I(world, s, a) <= I_threshold]
    if not safe:
        return 0
    x, y = s
    b = 1 if x < 5 else (0 if y > 0 else 1)
    return b if b in safe else safe[0]

# Results:
#   NRMO-like achieves Pr(endo_ruin <= 40) <= 0.02 cert.
#   Baseline policy does not satisfy the certificate.
\end{lstlisting}

\paragraph{Result interpretation.}

\begin{itemize}
\item NRMO-like policy with $I_{\mathrm{threshold}} = 0.55$ achieves
$\Pr(\mathrm{endogenous\_ruin} \le 40) \le 0.02$ certificate.
\item Baseline (unconstrained greedy) policy does \textbf{not}
satisfy the certificate.
\end{itemize}

% =====================================================================
\section*{Connections to Other Parts}

\begin{itemize}
\item These toy simulations are non-normative illustrations. The
canonical engineering implementation is in Part~X
(StrongEngine $\Omega$ Full v5.0, full civilisation simulator).
\item The irreversibility proxy $I(s, a)$ in Appendix~C.2 is the
operational analogue of the $\widehat{p}_R(a, b)$ ruin upper
bound in APCSO (Equation~\ref{eq:apcso-ruin-upper}).
\item The bounded-horizon certificate $\Pr(\tau_F \le T) \le
\varepsilon$ corresponds to the safe-policy class
$\mathrm{safe}(T, \varepsilon)$ in Part~V
(Definition~\ref{def:safe-policy}) and Part~VII.
\end{itemize}
